\chapter{Simplicial Complexes}
\label{ch:viii15-simplicial-complexes}

Simplicial complexes replace a topological space by vertices, edges, triangles,
tetrahedra, and higher-dimensional simplices glued face to face.  They are the
first combinatorial models used systematically in homology.  This chapter
develops abstract and geometric simplicial complexes, realizations, links and
stars, subdivisions, orientations, and simplicial maps, while leaving the full
chain-complex formalism to VIII/16.

\section{Geometric simplices}

Let
\[
v_0,\dots,v_k\in\mathbb R^N
\]
be affinely independent.

\begin{definition}[Geometric simplex]
\label{def:viii15-geometric-simplex}
The \(k\)-simplex spanned by these vertices is
\[
[v_0,\dots,v_k]
=
\left\{
\sum_{i=0}^k t_iv_i:
t_i\ge0,\ \sum_{i=0}^k t_i=1
\right\}.
\]
\end{definition}

The coefficients \(t_i\) are barycentric coordinates.

\section{Faces}

A face of
\[
[v_0,\dots,v_k]
\]
is obtained by retaining a subset of its vertices.

The \(i\)th codimension-one face is
\[
[v_0,\dots,\widehat v_i,\dots,v_k].
\]

\begin{definition}[Dimension]
\label{def:viii15-simplex-dim}
A simplex with \(k+1\) vertices has dimension \(k\).
\end{definition}

Thus vertices are \(0\)-simplices, edges are \(1\)-simplices, triangles are
\(2\)-simplices, and tetrahedra are \(3\)-simplices.

\section{Geometric simplicial complexes}

\begin{definition}[Geometric simplicial complex]
\label{def:viii15-geometric-complex}
A geometric simplicial complex \(K\) is a collection of geometric simplices
such that:
\begin{enumerate}
\item every face of a simplex in \(K\) also lies in \(K\);
\item the intersection of two simplices of \(K\) is either empty or a common
face of both.
\end{enumerate}
\end{definition}

The second condition rules out uncontrolled crossings.

\section{Abstract simplicial complexes}

\begin{definition}[Abstract simplicial complex]
\label{def:viii15-abstract-complex}
An abstract simplicial complex \(K\) on a vertex set \(V\) is a collection of
finite nonempty subsets of \(V\) such that
\[
\sigma\in K,\quad \varnothing\ne\tau\subset\sigma
\quad\Longrightarrow\quad
\tau\in K.
\]
\end{definition}

Members \(\sigma\in K\) are simplices, and
\[
\dim\sigma=|\sigma|-1.
\]

\section{Dimension of a complex}

\begin{definition}[Complex dimension]
\label{def:viii15-complex-dim}
\[
\dim K=\sup_{\sigma\in K}\dim\sigma.
\]
\end{definition}

The dimension may be finite or infinite.

\section{Subcomplexes}

\begin{definition}[Subcomplex]
\label{def:viii15-subcomplex}
A subcollection \(L\subset K\) is a subcomplex if it is itself a simplicial
complex: whenever \(\sigma\in L\), every nonempty face of \(\sigma\) also lies
in \(L\).
\end{definition}

\section{Geometric realization}

An abstract complex becomes a topological space by assigning basis vectors to
vertices and realizing each abstract simplex as their convex hull.

\begin{definition}[Geometric realization]
\label{def:viii15-realization}
The geometric realization of \(K\), denoted
\[
|K|,
\]
is the union of its realized simplices with the natural weak topology.
\end{definition}

Different geometric embeddings of the same abstract complex yield homeomorphic
realizations.

\section{Basic examples}

\begin{example}[Simplex and boundary]
\label{ex:viii15-simplex-boundary}
The full simplex on \(n+1\) vertices realizes as \(D^n\), while its proper-face
subcomplex realizes as
\[
S^{n-1}.
\]
\end{example}

\begin{example}[Cycle]
A polygonal cycle with \(m\ge3\) vertices realizes as \(S^1\).
\end{example}

\begin{example}[Triangulated disk]
Two triangles sharing one edge can realize a quadrilateral disk.
\end{example}

\section{Triangulations}

\begin{definition}[Triangulation]
\label{def:viii15-triangulation}
A triangulation of a topological space \(X\) is a simplicial complex \(K\)
together with a homeomorphism
\[
|K|\cong X.
\]
\end{definition}

Not every arbitrary decomposition is a triangulation; the simplicial
intersection conditions must hold.

\section{Stars}

\begin{definition}[Open star]
\label{def:viii15-open-star}
For a vertex \(v\in K\), the open star \(\operatorname{st}^\circ(v)\) is the
union of interiors of all simplices containing \(v\).
\end{definition}

\begin{definition}[Closed star]
\label{def:viii15-closed-star}
The closed star \(\operatorname{st}(v)\) is the subcomplex generated by all
simplices containing \(v\).
\end{definition}

The closed star is contractible: it is a cone with vertex \(v\).

\section{Links}

\begin{definition}[Link]
\label{def:viii15-link}
The link of a vertex \(v\) is
\[
\operatorname{lk}(v)
=
\{\sigma\in K:
v\notin\sigma,\ \sigma\cup\{v\}\in K\}.
\]
\end{definition}

Informally, the link records the directions surrounding \(v\).

\begin{proposition}[Star as a cone]
\label{prop:viii15-star-cone}
\[
\operatorname{st}(v)
=
v*\operatorname{lk}(v),
\]
the simplicial cone on the link.
\end{proposition}

\section{Manifold links}

In a triangulated \(n\)-manifold, an interior vertex should have link
homeomorphic to
\[
S^{n-1},
\]
while a boundary vertex has a link resembling
\[
D^{n-1}.
\]

This local criterion motivates later combinatorial manifold theory.

\section{Barycentric subdivision}

\begin{definition}[Barycentric subdivision]
\label{def:viii15-barycentric}
The barycentric subdivision \(\operatorname{sd}K\) has one vertex for each
nonempty simplex of \(K\).  A collection
\[
\sigma_0,\dots,\sigma_r
\]
forms a simplex in \(\operatorname{sd}K\) exactly when, after reordering,
\[
\sigma_0\subsetneq\sigma_1\subsetneq\cdots\subsetneq\sigma_r.
\]
\end{definition}

\begin{theorem}[Subdivision preserves realization]
\label{thm:viii15-subdivision-homeomorphism}
\[
|\operatorname{sd}K|\cong|K|.
\]
\end{theorem}

Subdivision refines the combinatorics without changing the underlying
topological space.

\section{Orientation of simplices}

An orientation of a \(k\)-simplex is an ordering
\[
[v_0,\dots,v_k]
\]
of its vertices, modulo even permutations.

Reversing orientation changes the sign:
\[
[v_{\pi(0)},\dots,v_{\pi(k)}]
=
\operatorname{sgn}(\pi)[v_0,\dots,v_k].
\]

This sign convention prepares the simplicial boundary operator in VIII/16.

\section{Boundary preview}

For an oriented simplex, the formal alternating boundary is
\[
\partial[v_0,\dots,v_k]
=
\sum_{i=0}^k
(-1)^i
[v_0,\dots,\widehat v_i,\dots,v_k].
\]

The central identity
\[
\partial^2=0
\]
will become the defining algebraic condition for chain complexes in VIII/16.

\section{Simplicial maps}

\begin{definition}[Simplicial map]
\label{def:viii15-simplicial-map}
Let \(K,L\) be abstract simplicial complexes.  A map on vertices
\[
f:V(K)\to V(L)
\]
is simplicial if for every simplex
\[
\sigma=\{v_0,\dots,v_k\}\in K,
\]
the image set
\[
f(\sigma)=\{f(v_0),\dots,f(v_k)\}
\]
is a simplex of \(L\).
\end{definition}

Distinct source vertices are allowed to have the same image.

\section{Realization of a simplicial map}

A simplicial map extends linearly over each simplex:
\[
|f|\left(\sum_i t_iv_i\right)
=
\sum_i t_i f(v_i).
\]

\begin{theorem}[Continuity of realization]
\label{thm:viii15-map-continuity}
A simplicial map induces a continuous map
\[
|f|:|K|\to|L|.
\]
\end{theorem}

\section{Composition}

\begin{proposition}[Functoriality]
\label{prop:viii15-map-functorial}
The composite of simplicial maps is simplicial, and
\[
|g\circ f|=|g|\circ|f|.
\]
\end{proposition}

\section{Simplicial isomorphisms}

\begin{definition}[Simplicial isomorphism]
\label{def:viii15-simplicial-iso}
A simplicial map is a simplicial isomorphism if it is bijective on vertices and
its inverse vertex map is also simplicial.
\end{definition}

A simplicial isomorphism induces a homeomorphism of realizations.

\section{Why subdivision helps maps}

A continuous map between polyhedra need not be simplicial with respect to a
given triangulation.  After sufficiently fine subdivisions, it can often be
approximated by a simplicial map.  This is the content of the simplicial
approximation theorem, developed later as needed.

\section{Bridge to chain complexes}

Once each oriented \(k\)-simplex is treated as a generator of an abelian group,
the alternating face formula becomes a homomorphism
\[
\partial_k:C_k\to C_{k-1}.
\]
The cancellation
\[
\partial_{k-1}\partial_k=0
\]
creates the chain-complex structure of VIII/16.


\section{Exercises}

\begin{exercise}\label{exr:viii15-01}
Define a geometric \(k\)-simplex.
\end{exercise}
\begin{hint}
Use convex combinations of \(k+1\) affinely independent vertices.

% VIII-PEDAGOGY-HINT-v1.1 exr:viii15-01 append
Require nonnegative barycentric coordinates summing to one; affine independence ensures the vertices span dimension \(k\) and coordinates are unique.
\end{hint}
\begin{solution}
\([v_0,\dots,v_k]=\{\sum t_iv_i:t_i\ge0,\sum t_i=1\}\).
\end{solution}

\begin{exercise}\label{exr:viii15-02}
How many vertices does a \(k\)-simplex have?
\end{exercise}
\begin{hint}
Dimension is one less than vertex count.

% VIII-PEDAGOGY-HINT-v1.1 exr:viii15-02 append
Check the cases of a vertex, edge, and triangle before expressing the general relation between affine dimension and vertex count.
\end{hint}
\begin{solution}
\(k+1\).
\end{solution}

\begin{exercise}\label{exr:viii15-03}
What is a face of a simplex?
\end{exercise}
\begin{hint}
Retain a subset of vertices.

% VIII-PEDAGOGY-HINT-v1.1 exr:viii15-03 append
Set the coordinates of omitted vertices to zero; use the chapter's convention that only nonempty faces are included.
\end{hint}
\begin{solution}
The simplex spanned by a nonempty subset of the original vertices.
\end{solution}

\begin{exercise}\label{exr:viii15-04}
State the intersection condition in a geometric simplicial complex.
\end{exercise}
\begin{hint}
Two simplices meet cleanly.

% VIII-PEDAGOGY-HINT-v1.1 exr:viii15-04 append
An arbitrary overlap, such as crossing edges without a shared vertex, fails the condition; identify the common face in each allowed intersection.
\end{hint}
\begin{solution}
Their intersection is empty or a common face of both.
\end{solution}

\begin{exercise}\label{exr:viii15-05}
Define an abstract simplicial complex.
\end{exercise}
\begin{hint}
Use finite subsets closed under nonempty subsets.

% VIII-PEDAGOGY-HINT-v1.1 exr:viii15-05 append
Test downward closure: including a triangle forces its edges and vertices to be present, while including its edges does not force the triangle.
\end{hint}
\begin{solution}
A family \(K\) of finite nonempty vertex sets such that every nonempty subset of a simplex is again a simplex.
\end{solution}

\begin{exercise}\label{exr:viii15-06}
What is \(\dim\sigma\) for an abstract simplex?
\end{exercise}
\begin{hint}
Count vertices.

% VIII-PEDAGOGY-HINT-v1.1 exr:viii15-06 append
Count elements of the vertex set, not its proper subsets; a singleton must have dimension zero.
\end{hint}
\begin{solution}
\(|\sigma|-1\).
\end{solution}

\begin{exercise}\label{exr:viii15-07}
Define a subcomplex.
\end{exercise}
\begin{hint}
Close under faces.

% VIII-PEDAGOGY-HINT-v1.1 exr:viii15-07 append
For each selected simplex, verify that every nonempty subset of its vertices is also selected and belongs to the original complex.
\end{hint}
\begin{solution}
A subcollection that contains every face of each of its simplices.
\end{solution}

\begin{exercise}\label{exr:viii15-08}
What is \(|K|\)?
\end{exercise}
\begin{hint}
Pass from combinatorics to topology.

% VIII-PEDAGOGY-HINT-v1.1 exr:viii15-08 append
Glue standard simplices along the common faces prescribed by the abstract complex and specify the resulting realization topology.
\end{hint}
\begin{solution}
The geometric realization of the abstract simplicial complex.
\end{solution}

\begin{exercise}\label{exr:viii15-09}
What space does a full \(n\)-simplex realize?
\end{exercise}
\begin{hint}
It is convex.

% VIII-PEDAGOGY-HINT-v1.1 exr:viii15-09 append
Convexity alone gives contractibility, not a ball homeomorphism; use the standard simplex's radial parametrization from an interior point.
\end{hint}
\begin{solution}
\(D^n\).
\end{solution}

\begin{exercise}\label{exr:viii15-10}
What does its boundary subcomplex realize?
\end{exercise}
\begin{hint}
Use the simplex boundary.

% VIII-PEDAGOGY-HINT-v1.1 exr:viii15-10 append
For \(n\ge1\), project radially from the barycenter to identify the simplex boundary with a sphere of dimension one less.
\end{hint}
\begin{solution}
\(S^{n-1}\).
\end{solution}

\begin{exercise}\label{exr:viii15-11}
Define a triangulation of \(X\).
\end{exercise}
\begin{hint}
Use a simplicial complex plus homeomorphism.

% VIII-PEDAGOGY-HINT-v1.1 exr:viii15-11 append
A homotopy equivalence is not enough: a triangulation includes a homeomorphism from the realization onto the given space.
\end{hint}
\begin{solution}
A complex \(K\) with a homeomorphism \(|K|\cong X\).
\end{solution}

\begin{exercise}\label{exr:viii15-12}
Define the closed star of a vertex.
\end{exercise}
\begin{hint}
Collect simplices containing the vertex and all their faces.

% VIII-PEDAGOGY-HINT-v1.1 exr:viii15-12 append
Include faces opposite the chosen vertex as well as simplices containing it; otherwise the proposed collection need not be a subcomplex.
\end{hint}
\begin{solution}
The subcomplex generated by all simplices containing the vertex.
\end{solution}

\begin{exercise}\label{exr:viii15-13}
Why is a closed star contractible?
\end{exercise}
\begin{hint}
It is a cone.

% VIII-PEDAGOGY-HINT-v1.1 exr:viii15-13 append
Join every point of the closed star to \(v\) within a containing simplex; barycentric interpolation gives a contraction that agrees on shared faces.
\end{hint}
\begin{solution}
\(\operatorname{st}(v)=v*\operatorname{lk}(v)\).
\end{solution}

\begin{exercise}\label{exr:viii15-14}
Define the link of \(v\).
\end{exercise}
\begin{hint}
Take faces opposite \(v\).

% VIII-PEDAGOGY-HINT-v1.1 exr:viii15-14 append
Impose both conditions: the simplex excludes \(v\), but adjoining \(v\) produces a simplex of \(K\).
\end{hint}
\begin{solution}
\(\operatorname{lk}(v)=\{\sigma:v\notin\sigma,\sigma\cup\{v\}\in K\}\).
\end{solution}

\begin{exercise}\label{exr:viii15-15}
What is the expected link of an interior vertex in an \(n\)-manifold triangulation?
\end{exercise}
\begin{hint}
One dimension lower.

% VIII-PEDAGOGY-HINT-v1.1 exr:viii15-15 append
In a combinatorial or PL manifold triangulation, examine a small vertex neighborhood as the cone on its link; do not silently extend this sphere-link criterion to all topological triangulations.
\end{hint}
\begin{solution}
\(S^{n-1}\).
\end{solution}

\begin{exercise}\label{exr:viii15-16}
What are vertices of \(\operatorname{sd}K\)?
\end{exercise}
\begin{hint}
Use simplices of \(K\).

% VIII-PEDAGOGY-HINT-v1.1 exr:viii15-16 append
Each nonempty original simplex contributes one new vertex at its barycenter, including original vertices viewed as zero-dimensional simplices.
\end{hint}
\begin{solution}
Nonempty simplices of \(K\), represented by their barycenters.
\end{solution}

\begin{exercise}\label{exr:viii15-17}
When do vertices form a simplex in \(\operatorname{sd}K\)?
\end{exercise}
\begin{hint}
Require a strict inclusion chain.

% VIII-PEDAGOGY-HINT-v1.1 exr:viii15-17 append
Check comparability under inclusion, not merely nonempty intersection; a simplex in the subdivision encodes a strict flag of original faces.
\end{hint}
\begin{solution}
When the corresponding simplices of \(K\) form a chain under inclusion.
\end{solution}

\begin{exercise}\label{exr:viii15-18}
What happens to the realization under barycentric subdivision?
\end{exercise}
\begin{hint}
Topology is unchanged.

% VIII-PEDAGOGY-HINT-v1.1 exr:viii15-18 append
Map a new vertex to the barycenter of its original simplex and extend affinely on flags; the pieces triangulate each original simplex without changing it.
\end{hint}
\begin{solution}
\(|\operatorname{sd}K|\cong|K|\).
\end{solution}

\begin{exercise}\label{exr:viii15-19}
Define an orientation of a simplex.
\end{exercise}
\begin{hint}
Order vertices modulo even permutations.

% VIII-PEDAGOGY-HINT-v1.1 exr:viii15-19 append
Compare the parity of two vertex orderings; one transposition reverses orientation while an even permutation preserves it.
\end{hint}
\begin{solution}
An equivalence class of vertex orderings under even permutations.
\end{solution}

\begin{exercise}\label{exr:viii15-20}
Write the alternating boundary of \([v_0,\dots,v_k]\).
\end{exercise}
\begin{hint}
Delete one vertex at a time.

% VIII-PEDAGOGY-HINT-v1.1 exr:viii15-20 append
Keep the surviving vertices in their inherited order and assign the sign from the position of the deleted vertex, starting at index zero.
\end{hint}
\begin{solution}
\(\partial[v_0,\dots,v_k]=\sum_i(-1)^i[v_0,\dots,\widehat v_i,\dots,v_k]\).
\end{solution}

\begin{exercise}\label{exr:viii15-21}
Define a simplicial map.
\end{exercise}
\begin{hint}
Images of vertex sets must remain simplices.

% VIII-PEDAGOGY-HINT-v1.1 exr:viii15-21 append
The map is specified on vertices, but its defining condition must hold for every simplex, not only for its edges.
\end{hint}
\begin{solution}
A vertex map \(f:V(K)\to V(L)\) sending every simplex of \(K\) to a simplex of \(L\).
\end{solution}

\begin{exercise}\label{exr:viii15-22}
May a simplicial map identify two vertices?
\end{exercise}
\begin{hint}
Injectivity is not required.

% VIII-PEDAGOGY-HINT-v1.1 exr:viii15-22 append
Test the collapse of one edge to a vertex; the image has lower dimension but remains a simplex, so simpliciality need not imply injectivity.
\end{hint}
\begin{solution}
Yes, provided the resulting image set is a simplex.
\end{solution}

\begin{exercise}\label{exr:viii15-23}
How is \(|f|\) defined on barycentric coordinates?
\end{exercise}
\begin{hint}
Extend linearly.

% VIII-PEDAGOGY-HINT-v1.1 exr:viii15-23 append
When several vertices have the same image, add their barycentric coefficients; nonnegativity and the total sum of one are preserved.
\end{hint}
\begin{solution}
\(|f|(\sum t_iv_i)=\sum t_if(v_i)\).
\end{solution}

\begin{exercise}\label{exr:viii15-24}
What is the bridge to VIII/16?
\end{exercise}
\begin{hint}
Turn oriented simplices into generators.

% VIII-PEDAGOGY-HINT-v1.1 exr:viii15-24 append
Form a free abelian group in each dimension, impose the orientation sign relation, and check cancellation of each codimension-two face in the double boundary.
\end{hint}
\begin{solution}
The alternating face maps become boundary homomorphisms in a chain complex.
\end{solution}


\section{Solved problem dossiers}

\begin{problem}[Triangle faces]\label{prob:viii15-01}
List all nonempty faces of the abstract simplex \(\{0,1,2\}\).
\end{problem}
\begin{solution}
The three vertices, the three edges \(\{0,1\},\{0,2\},\{1,2\}\), and the triangle \(\{0,1,2\}\).
\end{solution}

\begin{problem}[Tetrahedron count]\label{prob:viii15-02}
Count the \(0\)-, \(1\)-, \(2\)-, and \(3\)-faces of a tetrahedron.
\end{problem}
\begin{solution}
There are \(4\) vertices, \(\binom42=6\) edges, \(\binom43=4\) triangular faces, and one tetrahedron.
\end{solution}

\begin{problem}[Boundary of a triangle]\label{prob:viii15-03}
Identify the realization of the proper-face subcomplex of a 2-simplex.
\end{problem}
\begin{solution}
It is the three-edge cycle, homeomorphic to \(S^1\).
\end{solution}

\begin{problem}[Boundary of a tetrahedron]\label{prob:viii15-04}
What space is realized by the boundary complex of a 3-simplex?
\end{problem}
\begin{solution}
The four triangular faces glue to a 2-sphere \(S^2\).
\end{solution}

\begin{problem}[Invalid geometric complex]\label{prob:viii15-05}
Why do two triangles crossing through their interiors fail to form a geometric simplicial complex?
\end{problem}
\begin{solution}
Their intersection is not empty and is not a common face of the two triangles, violating the intersection axiom.
\end{solution}

\begin{problem}[Cycle complex]\label{prob:viii15-06}
Build an abstract simplicial complex realizing a square circle.
\end{problem}
\begin{solution}
Use vertices \(1,2,3,4\), all singleton vertices, and edges \(\{1,2\},\{2,3\},\{3,4\},\{4,1\}\). Its realization is \(S^1\).
\end{solution}

\begin{problem}[Filled versus unfilled triangle]\label{prob:viii15-07}
Compare the complexes with and without the 2-simplex \(\{0,1,2\}\).
\end{problem}
\begin{solution}
With the 2-simplex the realization is a disk; with only its proper faces the realization is a circle.
\end{solution}

\begin{problem}[Star and link in a triangle]\label{prob:viii15-08}
For vertex \(0\) in a filled triangle, identify its closed star and link.
\end{problem}
\begin{solution}
The star is the entire triangle and all faces incident/generated by it. The link consists of vertices \(1,2\) and edge \(\{1,2\}\), hence realizes as an interval.
\end{solution}

\begin{problem}[Link in a cycle]\label{prob:viii15-09}
What is the link of a vertex in a triangulated circle?
\end{problem}
\begin{solution}
It consists of the two neighboring vertices as two isolated \(0\)-simplices, realizing \(S^0\).
\end{solution}

\begin{problem}[Link in a triangulated surface]\label{prob:viii15-10}
What local link should an interior vertex have?
\end{problem}
\begin{solution}
A circle \(S^1\), matching the \(n=2\) manifold link rule.
\end{solution}

\begin{problem}[First barycentric subdivision of an edge]\label{prob:viii15-11}
Describe \(\operatorname{sd}\) of a single edge.
\end{problem}
\begin{solution}
The original two vertices and the edge barycenter become subdivision vertices; the edge is replaced by two edges meeting at the barycenter.
\end{solution}

\begin{problem}[Subdivision of a triangle]\label{prob:viii15-12}
How many small triangles appear in the first barycentric subdivision of one 2-simplex?
\end{problem}
\begin{solution}
There are \(3!=6\) maximal chains vertex \(\subset\) edge \(\subset\) triangle, hence six small triangles.
\end{solution}

\begin{problem}[Orientation reversal]\label{prob:viii15-13}
Compare \([v_0,v_1,v_2]\) and \([v_1,v_0,v_2]\).
\end{problem}
\begin{solution}
The transposition is odd, so the orientations differ by a minus sign.
\end{solution}

\begin{problem}[Triangle boundary]\label{prob:viii15-14}
Compute the formal oriented boundary of \([v_0,v_1,v_2]\).
\end{problem}
\begin{solution}
\(\partial[v_0,v_1,v_2]=[v_1,v_2]-[v_0,v_2]+[v_0,v_1]\).
\end{solution}

\begin{problem}[Boundary-squared preview]\label{prob:viii15-15}
Check formally that \(\partial^2[v_0,v_1,v_2]=0\).
\end{problem}
\begin{solution}
Apply \(\partial[v_i,v_j]=[v_j]-[v_i]\). Every vertex occurs twice with opposite signs, so the sum is zero.
\end{solution}

\begin{problem}[Constant simplicial map]\label{prob:viii15-16}
When is a constant vertex map \(K\to L\) simplicial?
\end{problem}
\begin{solution}
If its common image vertex belongs to \(L\), every simplex image is that singleton vertex, hence a simplex. Thus any constant map to a vertex is simplicial.
\end{solution}

\begin{problem}[Vertex identification]\label{prob:viii15-17}
Map a triangle's vertices \(0,1,2\) to \(a,a,b\). When is the map simplicial?
\end{problem}
\begin{solution}
It is simplicial provided \(\{a,b\}\) is a simplex of the target; the triangle image set is then the edge \(\{a,b\}\).
\end{solution}

\begin{problem}[Composition]\label{prob:viii15-18}
Prove a composite of simplicial maps is simplicial.
\end{problem}
\begin{solution}
If \(\sigma\in K\), then \(f(\sigma)\) is a simplex of \(L\), and \(g(f(\sigma))\) is a simplex of \(M\). Hence \(g\circ f\) is simplicial.
\end{solution}

\begin{problem}[Realization continuity]\label{prob:viii15-19}
Why is the realization of a simplicial map continuous?
\end{problem}
\begin{solution}
On each simplex it is affine and hence continuous; the formulas agree on common faces, and the weak topology of the realization glues them into a global continuous map.
\end{solution}

\begin{problem}[Simplicial isomorphism]\label{prob:viii15-20}
Why does a simplicial isomorphism induce a homeomorphism?
\end{problem}
\begin{solution}
Both the map and its inverse are simplicial, so their realizations are continuous and inverse to one another.
\end{solution}

\begin{problem}[Graph as a simplicial complex]\label{prob:viii15-21}
What restriction must be handled when viewing an arbitrary graph as an abstract simplicial complex?
\end{problem}
\begin{solution}
Abstract simplicial complexes do not allow loop edges with repeated vertices or multiple distinct edges with the same endpoint set. Subdivision can remove these obstacles.
\end{solution}

\begin{problem}[Triangulated disk Euler preview]\label{prob:viii15-22}
For a square cut by one diagonal, count vertices, edges, and triangles.
\end{problem}
\begin{solution}
There are \(4\) vertices, \(5\) edges, and \(2\) triangles. The alternating count is \(4-5+2=1\), foreshadowing disk Euler characteristic.
\end{solution}

\begin{problem}[Triangulated sphere Euler preview]\label{prob:viii15-23}
For the tetrahedral sphere, compute \(V-E+F\).
\end{problem}
\begin{solution}
\(4-6+4=2\), foreshadowing \(\chi(S^2)=2\).
\end{solution}

\begin{problem}[Bridge to chains]\label{prob:viii15-24}
Explain why orientations and alternating faces are exactly the data needed for simplicial homology.
\end{problem}
\begin{solution}
Oriented \(k\)-simplices become generators of free abelian groups, and the signed face sum defines the boundary map. The identity \(\partial^2=0\) makes these groups into the chain complex developed next.
\end{solution}
